\chapter{Descripción General y Arquitectura del Sistema}

\section{Descripción General del Proyecto}
El proyecto **LaptopRec Híbrido** es una plataforma integral de recomendación de computadoras portátiles diseñada bajo los principios de desacoplamiento de software (REST API) e hibridación en dos niveles. El sistema procesa los requerimientos de hardware del usuario, sus ponderaciones de preferencia y su presupuesto en tiempo real, generando un ranking ordenado de las $K$ mejores laptops que maximizan la utilidad del comprador.

\section{Pipeline de Procesamiento de Datos y EDA}
El flujo de datos del sistema sigue un procesamiento secuencial en pipeline desde la ingesta de catálogo hasta la inferencia híbrida, como se ilustra en la Figura \ref{fig:pipeline_datos}.

\begin{figure}[H]
    \centering
    \setlength{\unitlength}{1cm}
    \begin{picture}(14,6)
        % Cajas del Pipeline
        \put(0.5,4.5){\framebox(3.8,1.2){\shortstack{\textbf{1. Ingesta y EDA}\\[2pt]\small 150 Laptops / 1999 Ratings}}}
        \put(5.1,4.5){\vector(1,0){0.8}}
        
        \put(6.0,4.5){\framebox(3.8,1.2){\shortstack{\textbf{2. Normalización}\\[2pt]\small MinMax / One-Hot Encoding}}}
        \put(10.0,4.5){\vector(1,0){0.8}}
        
        \put(11.0,4.5){\framebox(2.8,1.2){\shortstack{\textbf{3. SVD / MAUT}\\[2pt]\small Vectorización TF-IDF}}}
        
        \put(2.4,4.5){\vector(0,-1){1.2}}
        
        \put(0.5,2.0){\framebox(3.8,1.3){\shortstack{\textbf{4. Nivel 1: Reglas}\\[2pt]\small Máscara Binaria $M_{rules}$}}}
        \put(5.1,2.6){\vector(1,0){0.8}}
        
        \put(6.0,2.0){\framebox(3.8,1.3){\shortstack{\textbf{5. Nivel 2: Ensamble}\\[2pt]\small $S_{hybrid} = \alpha U + \beta r + \gamma S$}}}
        \put(10.0,2.6){\vector(1,0){0.8}}
        
        \put(11.0,2.0){\framebox(2.8,1.3){\shortstack{\textbf{6. REST API / UI}\\[2pt]\small Top-K / 1vs1 / Chart.js}}}
    \end{picture}
    \caption{Diagrama del Pipeline de Procesamiento de Datos del Recomendador Híbrido.}
    \label{fig:pipeline_datos}
\end{figure}

\subsection{Análisis Exploratorio de Datos (EDA) y Preprocesamiento}
El conjunto de datos comprende 150 computadoras portátiles distribuidas en 6 marcas principales (ASUS, Apple, Dell, Lenovo, HP, MSI) y 1,999 calificaciones generadas por 100 usuarios sintéticos estructurados en 3 perfiles de comportamiento:
\begin{itemize}
    \item \textbf{Perfil Gamer / Power User}: Preferencia por GPUs dedicadas, VRAM $\ge 8$ GB y procesadores multi-núcleo ($>8$ núcleos).
    \item \textbf{Perfil Estudiante / Oficina}: Sensible al precio, preferencia por peso liviano ($\le 1.5$ kg) e integrada autonomía de batería.
    \item \textbf{Perfil Científico de Datos / IA}: Exigencia estricta de soporte NVIDIA CUDA, memoria RAM $\ge 16$ GB y resolución de pantalla elevada.
\end{itemize}

El preprocesamiento aplica escalamiento MinMax a las variables continuas (precio, RAM, cores, VRAM, peso) para transformar los atributos al dominio estandarizado $[0, 1]$.

\section{Arquitectura del Sistema en 2 Niveles}
La arquitectura del Recomendador Híbrido está dividida en dos capas funcionales superpuestas:

\subsection{Nivel 1: Filtro Estricto por Reglas Lógicas de Dominio}
Calcula una máscara binaria $M_{rules}(i) \in \{0, 1\}$ para cada laptop $i$ en el catálogo. Si la laptop excede el presupuesto máximo del usuario o no cumple con las restricciones mínimas del caso de uso seleccionado (ej. soporte CUDA para Deep Learning), $M_{rules}(i) = 0$, descartándola inmediatamente del ensamble.

\subsection{Nivel 2: Ensamble Ponderado Dinámico}
Para los elementos que superan el Nivel 1 ($M_{rules}(i) = 1$), se calcula el puntaje híbrido final $Score_{hybrid}(u, i)$ combinando tres modelos independientes:
\begin{equation}
Score_{hybrid}(u, i) = M_{rules}(i) \cdot \left[ \alpha \cdot U_{MAUT}(i) + \beta \cdot \hat{r}_{SVD\_norm}(u, i) + \gamma \cdot S_{content}(u, i) \right]
\end{equation}

Donde $\alpha, \beta, \gamma \ge 0$ son los pesos de ponderación que satisfacen $\alpha + \beta + \gamma = 1$.

\section{Arquitectura de Software y Módulos de Integración}
El sistema se implementa siguiendo el patrón de arquitectura limpia:
\begin{itemize}
    \item \textbf{Backend REST (Python FastAPI)}: Expone endpoints `/api/recommend`, `/api/laptops`, `/api/users` y `/api/metrics`.
    \item \textbf{Frontend SPA (HTML5, CSS3, JavaScript ES6+)}: Dashboard dinámico responsivo con soporte para selector multimoneda (USD, PEN, EUR), comparador cara a cara (1vs1) y renderizado de gráficos de tendencias mediante **Chart.js**.
\end{itemize}
